\scanpage{81}
\uline{Lemme} 1\qquad Soit
\[1\longrightarrow L\longrightarrow\widetilde\pi\;\text{---}\;\pi\;\text{---}\;1\]
une suite exacte de groupes, soit $(\unsure{x_i})_{i\in I}$ un syst. de générateurs de $\pi$,
\[(\unsure{a_j})_{j\in J}\qquad\text{\rule{5cm}{.4pt}}\qquad L,\]
$R_\alpha((x_i))=1$ $(\alpha\in A)$ un syst. de relations \unsure{définissantes} de $\pi$, $S_\beta((a_j))=1$ $(\beta\in B)$ un syst. de relations \unsure{définissantes} de $L$. Relevons les $x_i$ en des $\widetilde x_i\in\widetilde\pi$, et choisissons des mots $M_\alpha$ et $M_{i_0j_0}$ $(\alpha\in A,\ \deleted{\unsure{i\in I}}\ (i,j)\in I\times J)$ en \unsure{les} lettres $a_j$, de telle façon qu'on ait
\[
\left\{\begin{aligned}
(1)\quad R_\alpha((\widetilde x_i))&=M_\alpha((a_j))&&\forall\alpha\in A\\
(2)\quad\widetilde x_{i_0}a_{j_0}\widetilde x_{i_0}^{-1}&=M_{i_0j_0}((a_j))&&\forall(i_0,j_0)\in I\times J
\end{aligned}\right.
\]
\margin{relations $\widetilde R_\alpha$ ``\unsure{héritées}'' des rel. de $\pi$}
\margin{relations $\unsure{r_{(i,j)}}$ de ``conjugaison''}
Alors les relations (1) (2), plus les :
\[(3)\qquad\left\{S_\beta((a_j))=1\qquad(\beta\in B)\right.\]
\margin{``relations du \unsure{noyau}'' $S_\beta$}
sont un syst. de relations \unsure{définissantes} pour $\unsure{E|\widetilde\pi})$ p.r. aux générateurs $(\widetilde x_i)$, $(a_j)$.

\uline{NB} Si $L$ est central dans $\unsure{\widehat\pi|\widetilde\pi}$, \deleted{\unsure{les relations}}\inserted{\unsure{il y a}} intérêt à prendre les $M_{i_0j_0}$ les mots vides, donc (2) devient
\[(2')\qquad\widetilde x_i a_j\widetilde x_i^{-1}=a_j\qquad\text{i.e.}\quad[\widetilde x_i,a_j]=1\]
pour $\forall(i,j)\in I\times J$
\mspage{80}
